\chapter{Green Nail Syndrome}

\section*{Definition and Overview}
Green nail syndrome is a nail infection caused by the bacterium \textit{Pseudomonas aeruginosa}, which produces a green pigment that stains the nail. It typically affects people whose nails are chronically exposed to water or moisture, such as cleaners, food handlers, swimmers, and health care workers, and it often develops under artificial nails or after nail injury. The nail turns green, usually in a patch or band, and may be accompanied by separation of the nail from the nail bed (onycholysis). Green nail syndrome is not serious and is easily treated by keeping the nail dry, trimming the separated nail, and using antibacterial or antimicrobial treatments. It is important to distinguish it from other causes of green nail discolouration.

\section*{Epidemiology}
Green nail syndrome is not rare, but its exact prevalence is unknown because many cases are mild and not reported. It is more common in people with occupational or recreational water exposure, including cleaners, bartenders, swimmers, and nurses, and in those who wear artificial nails, which trap moisture against the natural nail. It affects adults more than children and is slightly more common in women, partly because of the greater use of artificial nails and nail cosmetics. Cases are more frequent in humid climates and in people with conditions that cause the nail to separate from the nail bed, such as psoriasis.

\section*{Aetiology and Pathophysiology}
Green nail syndrome is caused by \textit{Pseudomonas aeruginosa}, a bacterium that thrives in moist environments.

\begin{itemize}
    \item \textbf{The organism:} \textit{Pseudomonas aeruginosa} is a common environmental bacterium found in water, soil, and wet surfaces
    \item \textbf{Opportunistic infection:} it colonises the space between the nail and the nail bed when moisture is trapped there
    \item \textbf{Green pigment:} the bacterium produces pyocyanin, a blue-green pigment that stains the nail
    \item \textbf{Predisposing factors:} chronic water exposure, artificial nails, nail trauma, onycholysis (nail separation), and conditions such as psoriasis and diabetes
    \item \textbf{Infection of surrounding tissue:} in most cases the infection is confined to the nail, but it can occasionally spread to the skin and soft tissues, especially in people with weakened immunity
    \item \textbf{Contamination risk:} the infection can spread to other nails and, in health care settings, to wounds
\end{itemize}

\section*{Clinical Features}
The appearance and symptoms of green nail syndrome are usually characteristic.

\begin{itemize}
    \item A green, blue-green, or black-green discolouration of the nail, often in a patch or band
    \item Separation of the nail from the nail bed (onycholysis), often at the tip or edge
    \item The nail may be brittle, and debris can collect under it
    \item Usually no pain, unless there is secondary infection
    \item An unpleasant odour in some cases
    \item The surrounding skin is usually normal, though it may become infected in severe cases
    \item One or several nails can be affected
\end{itemize}

\section*{Differential Diagnosis}
Green nail discolouration has other causes that should be considered.

\begin{itemize}
    \item Trauma-related bruising, which may appear dark but is not typically green
    \item Fungal nail infection (onychomycosis), which usually causes thickening and yellow-white discolouration
    \item Psoriasis of the nail, which causes pitting, ridging, and separation
    \item Bacterial infection with other organisms
    \item Nail pigmentation from medicines or systemic illness
    \item Melanoma of the nail bed, which is rare but important to consider when there is a new dark streak
    \item Staining from nail polish, hair dye, or occupational substances
    \item Melanonychia and other pigmentary changes
\end{itemize}

\section*{When to Seek Medical Care}
See a doctor or podiatrist if the green discolouration persists despite keeping the nail dry, if the nail is painful, if the surrounding skin becomes red and swollen, or if you have diabetes or a weakened immune system.

\textbf{Call 000 (or local emergency services) for:}
\begin{itemize}
    \item Spreading redness, warmth, pain, or swelling around the nail with fever, which may indicate a spreading infection (cellulitis)
    \item Signs of severe infection in a person with diabetes or a weakened immune system
    \item Sudden severe pain or swelling of a finger or toe
\end{itemize}

\section*{Diagnosis and Investigations}
The diagnosis is usually clinical, based on the typical appearance.

\begin{itemize}
    \item \textbf{Clinical examination:} the green staining and nail separation are usually diagnostic
    \item \textbf{Culture:} a swab of the nail or subungual debris can confirm \textit{Pseudomonas}
    \item \textbf{Microscopy:} a nail clipping examined under a microscope can exclude fungal infection
    \item \textbf{Wood's lamp:} some pseudomonas infections fluoresce under ultraviolet light
    \item \textbf{Assessment of risk factors:} water exposure, artificial nails, and underlying conditions
    \item \textbf{Referral:} a dermatologist or podiatrist may be involved in persistent or recurrent cases
\end{itemize}

\section*{Management}
Treatment focuses on drying the nail and eliminating the bacterium.

\begin{itemize}
    \item \textbf{Keep the nail dry:} the single most important measure; dry hands and feet thoroughly after washing, and avoid prolonged immersion
    \item \textbf{Remove artificial nails:} artificial nails trap moisture and must be removed
    \item \textbf{Trim the separated nail:} removing the detached portion allows the nail bed to dry and the nail to regrow normally
    \item \textbf{Antimicrobial treatments:} antiseptic solutions such as dilute bleach soaks, acetic acid, or antibiotic drops applied to the nail bed
    \item \textbf{Topical antibiotics:} such as gentamicin or mupirocin in some cases
    \item \textbf{Oral antibiotics:} rarely needed, but used if the surrounding skin is infected or the infection is persistent
    \item \textbf{Protective measures:} wearing gloves for wet work and avoiding trauma to the nails
    \item \textbf{Follow-up:} the nail grows out over several months, and the green colour resolves as the infected portion is trimmed away
\end{itemize}

\section*{Complications}
\begin{itemize}
    \item Persistent or recurrent infection if the nail is not kept dry
    \item Spread of the infection to other nails
    \item Infection of the surrounding skin (paronychia or cellulitis)
    \item In people with weakened immunity or diabetes, more serious soft tissue infection
    \item Chronic nail separation and cosmetic concerns
    \item Secondary fungal infection of the separated nail
    \item Rarely, spread to bone in severely immunocompromised people
\end{itemize}

\section*{Prognosis and Outcomes}
Green nail syndrome is a benign condition with an excellent prognosis. It resolves once the nail is kept dry and the infected portion is trimmed, and the nail grows back normally over three to six months. Recurrence is common if the underlying moisture exposure continues, so prevention through drying, gloves, and removing artificial nails is the key to lasting resolution. Complications are uncommon in healthy people.

\section*{Prevention and Self-Care}
\begin{itemize}
    \item Dry hands and feet thoroughly after washing, swimming, or bathing, including under the nails
    \item Wear waterproof gloves for prolonged wet work, such as cleaning and food handling
    \item Avoid wearing artificial nails, which trap moisture
    \item Trim nails straight across and keep them short
    \item Avoid trauma to the nails and treat nail injuries promptly
    \item If a nail separates, keep the area clean and dry
    \item Use moisturiser to prevent dry skin, but keep nail folds dry
    \item See a health professional early if the nail is painful, the skin around it is red, or you have diabetes or a weakened immune system
\end{itemize}

\section*{Sources and Further Reading}
The following reputable sources provide further information for healthcare students and patients seeking reliable, current advice about green nail syndrome.

\begin{itemize}
    \item DermNet NZ. \textit{Green nail syndrome.} https://dermnetnz.org/
    \item Healthdirect Australia. \textit{Nail problems: green nail and other discolouration.} https://www.healthdirect.gov.au/
    \item Australasian College of Dermatologists. \textit{Nail disorders information.} https://www.dermcoll.edu.au/
    \item Australian Podiatry Association. \textit{Nail care and conditions.} https://www.podiatryaustralia.org.au/
    \item NHS. \textit{Nail problems: causes and treatment.} https://www.nhs.uk/
    \item StatPearls. \textit{Green nail syndrome.} https://www.ncbi.nlm.nih.gov/books/
\end{itemize}
